% GENERATED FILE. Edit canonical lessons, then run npm run generate:books.
% canonical-source-hash: fnv1a64:5db3af8227346465
% canonical-lessons: MW-C176-dhakno, MW-C176-badhno, MW-C176-gath, MW-C176-sakal, MW-C176-tagno

\chapter{Covering, Tying and Hanging}
\label{ch:mw-covering-tying-and-hanging}
\hlchaptermodality{176}

\begin{chapteropening}
\textbf{By the end of this chapter:} \emph{I can say that I cover the milk, tie a knot and hang things up.}

Complete the last lesson of chapter 176: I can say that I cover the milk, tie a knot and hang things up.
\end{chapteropening}

\section[\texorpdfstring{\mw{ढांकणो} (ḍhā̃kṇo)}{ढांकणो (ḍhā̃kṇo)}]{\mw{ढांकणो} (ḍhā̃kṇo) — to cover}

\label{lesson:MW-C176-dhakno}



\begin{quote}
Before the new one: say the Marwadi for a small water pot, then the Marwadi for a spouted water pot.
\end{quote}

\subsection*{You'll want to know: \mw{ढांकणो}}
\textbf{\mw{ढांकणो}} — \emph{ḍhā̃kṇo} — \textquotedblleft{}to cover\textquotedblright{}.

\begin{quote}
\mw{मां} \mw{दूध} \mw{ढांकै} \mw{है।}
\end{quote}

\begin{quote}
\emph{mā̃ dūdh ḍhā̃kai hai.}
\end{quote}

\begin{quote}
Mother covers the milk.
\end{quote}

Past (I …): \mw{म्हैं} \mw{दूध} \mw{ढांक्यो}, \emph{mhaĩ dūdh ḍhā̃kyo}.

\begin{cousinweb}
\textbf{\mw{ढांकणो}} (\emph{ḍhā̃kṇo}) is inherited from a reconstructed \emph{*ḍhakk}, \textquotedblleft{}cover\textquotedblright{}; CDIAL lists Marwari \emph{ḍhā̃kṇo} under it (CDIAL 5574).
\end{cousinweb}

\begin{grammarlens}[title={What you've built}]
One more word for this chapter.
\end{grammarlens}

\subsection*{Your turn}
\begin{itemize}
\raggedright
  \item \emph{Say it:} \emph{ḍhā̃kṇo}
  \item \emph{Say it:} \emph{ḍhā̃kṇo}, once more
  \item \emph{Say it:} \emph{mā̃ dūdh ḍhā̃kai hai.} — Mother covers the milk.
  \item \emph{Recall:} say the Marwadi for a small water pot, then the Marwadi for a spouted water pot, then say \emph{ḍhā̃kṇo} again
\end{itemize}

\begin{tcolorbox}[breakable,colback=teal!4,colframe=teal!35!black,title={Before you move on}]
What does \mw{ढांकणो} mean? (\textquotedblleft{}To cover\textquotedblright{}.) Say it once more.
\end{tcolorbox}
\section[\texorpdfstring{\mw{बांधणो} (bā̃dhṇo)}{बांधणो (bā̃dhṇo)}]{\mw{बांधणो} (bā̃dhṇo) — to tie}

\label{lesson:MW-C176-badhno}



\begin{quote}
Before the new one: say the Marwadi for a spouted water pot, then the Marwadi for to cover.
\end{quote}

\subsection*{You'll want to know: \mw{बांधणो}}
\textbf{\mw{बांधणो}} — \emph{bā̃dhṇo} — \textquotedblleft{}to tie\textquotedblright{}.

\begin{quote}
\mw{हूं} \mw{गाय} \mw{बांधूं} \mw{हूं।}
\end{quote}

\begin{quote}
\emph{hū̃ gāy bā̃dhū̃ hū̃.}
\end{quote}

\begin{quote}
I tie up the cow.
\end{quote}

Past (I …): \mw{म्हैं} \mw{गाय} \mw{बांधी}, \emph{mhaĩ gāy bā̃dhī}.

\begin{cousinweb}
\textbf{\mw{बांधणो}} (\emph{bā̃dhṇo}) is inherited from Sanskrit \emph{bandhati}, \textquotedblleft{}binds\textquotedblright{}; CDIAL lists the Old Marwari verb under it (CDIAL 9139).
\end{cousinweb}

\begin{grammarlens}[title={What you've built}]
One more word for this chapter.
\end{grammarlens}

\subsection*{Your turn}
\begin{itemize}
\raggedright
  \item \emph{Say it:} \emph{bā̃dhṇo}
  \item \emph{Say it:} \emph{bā̃dhṇo}, once more
  \item \emph{Say it:} \emph{hū̃ gāy bā̃dhū̃ hū̃.} — I tie up the cow.
  \item \emph{Recall:} say the Marwadi for a spouted water pot, then the Marwadi for to cover, then say \emph{bā̃dhṇo} again
\end{itemize}

\begin{tcolorbox}[breakable,colback=teal!4,colframe=teal!35!black,title={Before you move on}]
What does \mw{बांधणो} mean? (\textquotedblleft{}To tie\textquotedblright{}.) Say it once more.
\end{tcolorbox}
\section[\texorpdfstring{\mw{गांठ} (gā̃ṭh)}{गांठ (gā̃ṭh)}]{\mw{गांठ} (gā̃ṭh) — a knot}

\label{lesson:MW-C176-gath}



\begin{quote}
Before the new one: say the Marwadi for to cover, then the Marwadi for to tie.
\end{quote}

\subsection*{You'll want to know: \mw{गांठ}}
\textbf{\mw{गांठ}} — \emph{gā̃ṭh} — \textquotedblleft{}a knot\textquotedblright{}.

\begin{quote}
\mw{गांठ} \mw{घणी} \mw{कड़क} \mw{है।}
\end{quote}

\begin{quote}
\emph{gā̃ṭh ghaṇī kaṛak hai.}
\end{quote}

\begin{quote}
The knot is very tight.
\end{quote}

A feminine noun.

\begin{cousinweb}
\textbf{\mw{गांठ}} (\emph{gā̃ṭh}) is inherited from Sanskrit \emph{granthi}, \textquotedblleft{}a knot\textquotedblright{}; CDIAL lists Hindi and Gujarati \emph{gā̃ṭh} (CDIAL 4354). CDIAL gives no Marwari form in this entry.
\end{cousinweb}

\begin{grammarlens}[title={What you've built}]
One more word for this chapter.
\end{grammarlens}

\subsection*{Your turn}
\begin{itemize}
\raggedright
  \item \emph{Say it:} \emph{gā̃ṭh}
  \item \emph{Say it:} \emph{gā̃ṭh}, once more
  \item \emph{Say it:} \emph{gā̃ṭh ghaṇī kaṛak hai.} — The knot is very tight.
  \item \emph{Recall:} say the Marwadi for to cover, then the Marwadi for to tie, then say \emph{gā̃ṭh} again
\end{itemize}

\begin{tcolorbox}[breakable,colback=teal!4,colframe=teal!35!black,title={Before you move on}]
What does \mw{गांठ} mean? (\textquotedblleft{}A knot\textquotedblright{}.) Say it once more.
\end{tcolorbox}
\section[\texorpdfstring{\mw{सांकळ} (sā̃kaḷ)}{सांकळ (sā̃kaḷ)}]{\mw{सांकळ} (sā̃kaḷ) — a door chain}

\label{lesson:MW-C176-sakal}



\begin{quote}
Before the new one: say the Marwadi for to tie, then the Marwadi for a knot.
\end{quote}

\subsection*{You'll want to know: \mw{सांकळ}}
\textbf{\mw{सांकळ}} — \emph{sā̃kaḷ} — \textquotedblleft{}a door chain\textquotedblright{}.

\begin{quote}
\mw{किवाड़} \mw{माथै} \mw{सांकळ} \mw{है।}
\end{quote}

\begin{quote}
\emph{kivāṛ māthai sā̃kaḷ hai.}
\end{quote}

\begin{quote}
There's a chain on the door.
\end{quote}

A feminine noun.

\begin{cousinweb}
\textbf{\mw{सांकळ}} (\emph{sā̃kaḷ}) is inherited from Sanskrit \emph{śṛṅkhala}, \textquotedblleft{}a chain\textquotedblright{}; CDIAL lists Gujarati \emph{sā̃kaḷ} and Hindi \emph{sā̃kal} (CDIAL 12580). CDIAL gives no Marwari form in this entry.
\end{cousinweb}

\begin{grammarlens}[title={What you've built}]
One more word for this chapter.
\end{grammarlens}

\subsection*{Your turn}
\begin{itemize}
\raggedright
  \item \emph{Say it:} \emph{sā̃kaḷ}
  \item \emph{Say it:} \emph{sā̃kaḷ}, once more
  \item \emph{Say it:} \emph{kivāṛ māthai sā̃kaḷ hai.} — There's a chain on the door.
  \item \emph{Recall:} say the Marwadi for to tie, then the Marwadi for a knot, then say \emph{sā̃kaḷ} again
\end{itemize}

\begin{tcolorbox}[breakable,colback=teal!4,colframe=teal!35!black,title={Before you move on}]
What does \mw{सांकळ} mean? (\textquotedblleft{}A door chain\textquotedblright{}.) Say it once more.
\end{tcolorbox}
\section[\texorpdfstring{\mw{टांगणो} (ṭā̃gṇo)}{टांगणो (ṭā̃gṇo)}]{\mw{टांगणो} (ṭā̃gṇo) — to hang up}

\label{lesson:MW-C176-tagno}



\begin{quote}
Before the new one: say the Marwadi for a knot, then the Marwadi for a door chain.
\end{quote}

\subsection*{You'll want to know: \mw{टांगणो}}
\textbf{\mw{टांगणो}} — \emph{ṭā̃gṇo} — \textquotedblleft{}to hang up\textquotedblright{}.

\begin{quote}
\mw{हूं} \mw{कमीज} \mw{खूंटी} \mw{माथै} \mw{टांगूं} \mw{हूं।}
\end{quote}

\begin{quote}
\emph{hū̃ kamīj khū̃ṭī māthai ṭā̃gū̃ hū̃.}
\end{quote}

\begin{quote}
I hang my shirt on the peg.
\end{quote}

Past (I …): \mw{म्हैं} \mw{कमीज} \mw{टांगी}, \emph{mhaĩ kamīj ṭā̃gī}.

\begin{cousinweb}
\textbf{\mw{टांगणो}} (\emph{ṭā̃gṇo}) comes from a reconstructed \emph{*ṭaṅg}, \textquotedblleft{}hang up\textquotedblright{}; CDIAL lists Hindi \emph{ṭā̃gnā} and Gujarati \emph{ṭā̃gvũ} (CDIAL 5436). CDIAL gives no Marwari form in this entry.
\end{cousinweb}

\begin{grammarlens}[title={What you've built}]
One more word for this chapter.
\end{grammarlens}

\subsection*{Your turn}
\begin{itemize}
\raggedright
  \item \emph{Say it:} \emph{ṭā̃gṇo}
  \item \emph{Say it:} \emph{ṭā̃gṇo}, once more
  \item \emph{Say it:} \emph{hū̃ kamīj khū̃ṭī māthai ṭā̃gū̃ hū̃.} — I hang my shirt on the peg.
  \item \emph{Recall:} say the Marwadi for a knot, then the Marwadi for a door chain, then say \emph{ṭā̃gṇo} again
\end{itemize}

\begin{tcolorbox}[breakable,colback=teal!4,colframe=teal!35!black,title={Before you move on}]
What does \mw{टांगणो} mean? (\textquotedblleft{}To hang up\textquotedblright{}.) Say it once more.
\end{tcolorbox}
